\documentclass[letterpaper,12pt]{article}
\usepackage[spanish]{babel}
\usepackage[utf8]{inputenc}
\usepackage{graphicx}
\usepackage{background}
\usepackage{geometry}
\usepackage{eso-pic}
\usepackage{tikz}

% Ajusta los márgenes según tus preferencias
\geometry{
paperheight=25cm,
paperwidth=21cm, 
top=2cm, bottom=2cm, left=2.5cm, right=2.5cm}

\backgroundsetup{
  scale=4,
  angle=0,
  color=black,
  opacity=0.12,
  position=current page.south,
  hshift = -0.5cm,
  vshift=0.7cm,
  contents={%
\(\displaystyle \rho_{t} + (\rho v)_{x} = 0\)\hspace{6mm}
%\(\displaystyle (\rho v)_{t} + (\rho v^{2} + p)_{x} = 0\)\hspace{6mm}
%\(\displaystyle (\rho E)_{t} + [v(\rho E + p)]_{x} = 0\)
  }
}


\begin{document}
\pagestyle{empty}

\def \downleft{current page.south west}

\AddToShipoutPicture{%
  \begin{tikzpicture}[remember picture, overlay]
    \draw[line width=1pt] (1cm,1cm) rectangle (20cm,24cm);
  \end{tikzpicture}
}

\begin{center}
	\includegraphics[width=0.55\textwidth]{1.png}\hfill
    \includegraphics[width=0.28\textwidth]{usaclogo.png}
    

    \vspace{1cm}

    %\textbf{\LARGE Acto de Graduación}
    {\large El estudiante}

    \vspace{1cm}

    \textbf{\Large Rodrigo Rafael Castillo Chong}

    \vspace{1cm}
\large
    {tiene el honor de invitarle al solemne acto de graduación y examen público previo a serle conferido el título de}

    \vspace{1cm}

    \textbf{\Large Licenciado en Física Aplicada}

    \vspace{1cm}

    {\large que tendrá lugar en el \textbf{Auditorio Francisco Vela} de la facultad de Ingeniería, en la }

    \vspace{1cm}
    {\Large \textbf{Universidad de San Carlos de Guatemala}}
    \vspace{1cm}

    \textbf{\large [Miércoles 13 de marzo] \hspace{0.5cm} [3:00pm]}

    \vspace{1cm}

    \textbf{\large [Recepción: Vesuvio Majadas, 4:30pm]}

\end{center}

\end{document}
